% =====================================================================
% RASCUNHO (2026-11-06): texto da conclusão, para a autora trabalhar.
%   Os números de formato e de hardware devem ser conferidos com os Resultados quando
%   eles estiverem escritos. Trechos marcados com [completar] dependem deles.
% =====================================================================

\textcolor{red}{Este trabalho estudou o BCD-MMSE no \textit{uplink} de um sistema MIMO massivo com processamento descentralizado e ruído colorido. O foco foi a inversão da matriz $\mathbf G_c$ em cada DU, feita em ponto fixo em um FPGA. Um simulador em Python mediu o efeito do formato de ponto fixo nas curvas de SER. Depois, o módulo de inversão foi implementado em hardware, com uma arquitetura para cada condição de teste.}

\textcolor{red}{[completar] Formatos escolhidos para cada grupo de configurações e aumento de SNR de cada um, em relação ao ponto flutuante (Seção~\ref{sec:formato}).}

\textcolor{red}{[completar] Recursos e potência das arquiteturas. Todas atendem o limite de vazão da sua condição de teste.}

\textcolor{red}{As arquiteturas foram dimensionadas pela vazão, e não pelo atraso. Nas condições de baixa mobilidade, elas distribuem as inversões ao longo do tempo de coerência e gastam menos área e potência. Em troca, o atraso até a última inversão de um bloco fica maior. Parte do atraso vem do próprio método de estimação, pois juntar as $N$ amostras de ruído leva tempo.}

